\chapter{Intersections and sections}\label{ch:intersections}

Intersections are the bridge between free-form geometry and topology. Trimming
(\cref{ch:trimming}), sewing and Boolean operations (\cref{ch:solids}) and the slicing of
solids (\cref{ch:slicing}) are all built on the queries described in this chapter. \NF{}
offers four intersection problems behind a single overloaded function \apiref{intersect}:

\begin{center}
\begin{tabularx}{\linewidth}{@{}lXl@{}}
\toprule
Call & Problem & Result \\
\midrule
\code{intersect(a, b)} & curve / curve & \code{list[CurveCurveIntersection]} \\
\code{intersect(c, s)} & curve / surface & \code{list[CurveSurfaceIntersection]} \\
\code{intersect(c, plane)} & curve / infinite plane & \code{list[CurveSurfaceIntersection]} \\
\code{intersect(s1, s2[, options])} & surface / surface & \code{SurfaceIntersection} \\
\bottomrule
\end{tabularx}
\end{center}

Python selects the overload from the argument types: a \code{NurbsCurve} and a
\code{NurbsSurface} give the curve/surface variant, a \code{Plane} as second argument the
curve/plane variant. In addition, \apiref{planar_section} and \apiref{planar_sections} slice a
single \emph{surface} with planes, and \apiref{closest_point} /
\apiref{closest_point_local} project points onto curves and surfaces. Sections of whole
\emph{solids} (closed shells) are a separate topic and are covered in \cref{ch:slicing}.

All intersection functions are pure queries: they never modify their inputs, they may be
called concurrently from several threads on shared objects, and they report what they found
together with the residual gap, so that the caller can judge the quality of every result.

\section{Tolerances}\label{sec:isect-tolerances}

Every intersection routine takes a tolerance. Curve/curve, curve/surface and curve/plane
intersections take a \apiref{Tolerance} directly. The surface/surface intersection takes an
options object,
\apiref{SurfaceIntersectionOptions}, which contains a \code{Tolerance} plus two accuracy
parameters for the traced curves.

\begin{center}
\begin{tabularx}{\linewidth}{@{}llX@{}}
\toprule
Field & Default & Meaning \\
\midrule
\code{Tolerance.length} & $10^{-10}$ & Model-space distance below which two points count as
coincident. An intersection point is reported when the two objects come within this
distance. \\
\code{Tolerance.parametric} & $10^{-10}$ & Parameter differences below which two parameter
values are equal (merging of duplicates at Bézier piece boundaries). \\
\code{Tolerance.angular} & $10^{-9}$ & Angles (radians) below which directions are parallel. \\
\bottomrule
\end{tabularx}
\end{center}

The defaults are deliberately tight because \NF{} geometry is usually exact: lines, circles,
planes, cylinders and spheres are represented exactly, so their intersections can be computed
to rounding accuracy. The printed representation of a default \code{Tolerance()} object shows
\code{length=0.000000}, which is merely the six-digit formatting of $10^{-10}$.

\begin{tip}[When to increase the tolerance]
Increase \code{Tolerance.length} when the inputs themselves are approximations, e.g. fitted
measurement data, STEP imports from other systems, or curves produced by an earlier
intersection with a fit tolerance of $10^{-7}$. A tolerance smaller than the error of the
inputs makes tangential contacts disappear (the curves miss each other by more than the
tolerance) or splits a single crossing into several near-duplicate points.
\end{tip}

The tolerance model as a whole -- which operations are exact and which approximate, how
tolerances propagate -- is summarised in \cref{ch:concepts} and \cref{app:tolerances}.

\section{Curve/curve intersections}\label{sec:isect-cc}

\code{intersect(a, b, tolerance=Tolerance())} returns all points where the curves $a(s)$ and
$b(t)$ come within \code{tolerance.length} of each other, sorted by the parameter $s$ on the
first curve. Each result is a \apiref{CurveCurveIntersection} with four fields:

\begin{center}
\begin{tabularx}{\linewidth}{@{}lX@{}}
\toprule
Field & Meaning \\
\midrule
\code{s} & parameter on the first curve \code{a} \\
\code{t} & parameter on the second curve \code{b} \\
\code{point} & midpoint of $a(s)$ and $b(t)$ (a \code{Vec3}) \\
\code{distance} & the remaining gap $|a(s)-b(t)|$, at most \code{tolerance.length} \\
\bottomrule
\end{tabularx}
\end{center}

The curves are genuinely three-dimensional: two skew lines that cross in a projection but
miss each other in space have no intersection. The following example intersects a line with
a circle, a tangent line with the same circle, and two skew lines; it then turns to
curve/surface and curve/plane intersections, which are described in the next sections.

\example{09_curve_intersections.py}

Running the script prints

\begin{lstlisting}[style=text]
line/circle: s = 0.211325, t = 0.164730, point = (-1.732051, +1.000000), gap 0.0e+00
line/circle: s = 0.788675, t = 0.835270, point = (+1.732051, +1.000000), gap 5.0e-16
expected x = +-sqrt(3) = +-1.732051
tangent line: 1 contact point(s)
skew lines: 0 intersections
\end{lstlisting}

Note that the parameter $t$ on the circle is \emph{not} proportional to the angle: circles
and arcs are rational quadratic NURBS (\cref{ch:curves}), whose parametrisation is not
arc-length. Use \code{arc_length} or \code{closest_point} if you need angles or lengths.

\subsection{How it works}

The curves are decomposed into Bézier pieces, and pairs of pieces whose bounding boxes
overlap are subdivided recursively. A pair becomes a \emph{leaf} when it can contain at most
one intersection: the cones containing all tangent directions of the two pieces are
separated (also against the reversed orientation). A second intersection would require the
chord between the two points to lie in both cones, which is impossible. Within each leaf a
box-constrained Gauss--Newton iteration solves $a(s)-b(t)=0$, which converges quadratically
for transversal intersections. Candidates found in neighbouring leaves (at piece boundaries)
are merged when they are close or when the curve point at their parameter midpoint is still
within tolerance.

Tangential contacts never separate the tangent cones. There the subdivision stops at a size
limit ($10^{-6}$ of the model size) and the Gauss--Newton iteration finds the point of
closest approach; it is reported if the gap is within \code{tolerance.length}. That is why
the tangent line in the example yields exactly one contact point.

\begin{warning}[Overlapping curves]
Coincident (overlapping) curve portions -- two lines along the same segment, an arc lying on
a circle -- are not detected as overlaps. The result then consists of sample points along
the common part. Test for overlaps yourself, e.g. with \code{closest_point} at a few points
of one curve, before relying on the intersection points.
\end{warning}

\section{Curve/surface intersections}\label{sec:isect-cs}

The call \code{intersect(curve, surface)}, with an optional \code{tolerance} as third
argument, returns a list of intersection records, sorted by the curve parameter. Each one is a
\apiref{CurveSurfaceIntersection}:

\begin{center}
\begin{tabularx}{\linewidth}{@{}lX@{}}
\toprule
Field & Meaning \\
\midrule
\code{t} & parameter on the curve \\
\code{u}, \code{v} & parameters on the surface \\
\code{point} & midpoint of $c(t)$ and $S(u,v)$ \\
\code{distance} & the gap $|c(t)-S(u,v)|$, at most \code{tolerance.length} \\
\bottomrule
\end{tabularx}
\end{center}

The algorithm is the same as for two curves: subdivision of the curve and the surface into
Bézier pieces with bounding-box rejection, until the height of the curve piece over the
surface piece is monotone -- the angle between all curve tangents and all surface normals
stays below $90^\circ - (\alpha+\beta)$, where $\alpha,\beta$ are the opening angles of the
tangent and normal cones. Gauss--Newton on $c(t) - S(u,v) = 0$ then refines each candidate.

In the example above, a helix with three turns (pitch 1, radius 1) pierces a tilted
rectangular patch $z = 0.75\,(y+2)$:

\begin{lstlisting}[style=text]
helix/tilted plane: 5 points
  t = 0.750000 turns, (u, v) = (0.500000, 0.250000), z = 0.750000
  t = 0.818603 turns, (u, v) = (0.604455, 0.272868), z = 0.818603
  t = 1.500000 turns, (u, v) = (0.250000, 0.500000), z = 1.500000
  t = 2.181397 turns, (u, v) = (0.604455, 0.727132), z = 2.181397
  t = 2.250000 turns, (u, v) = (0.500000, 0.750000), z = 2.250000
\end{lstlisting}

The surface parameters $(u,v)$ refer to the patch's own parametrisation. For a patch made
with \code{make_rectangle(origin, edge_u, edge_v)} they are the fractions along the two
edges.

\begin{note}
Curve portions lying \emph{in} the surface (an edge lying in a plane) are not detected as
such by the plain curve/surface intersection. The Boolean operations handle this case
separately (\cref{sec:boolean-algorithm}); if you need it yourself, check a few points of the
curve with \code{closest_point}.
\end{note}

\section{Curve/plane intersections and the Plane type}\label{sec:isect-plane}

\apiref{Plane} is an infinite plane, \code{Plane(point, normal=(0, 0, 1))}. The normal is
normalised internally; \code{plane.point} and \code{plane.normal} return the defining data,
and \code{plane.distance(p)} the \emph{signed} distance of a point (positive on the side the
normal points to).

\code{intersect(curve, plane)} needs no surface patch: the curve is intersected with the
infinite plane. The results are again \code{CurveSurfaceIntersection} objects; their
\code{u}, \code{v} fields are the coordinates of the intersection point along two orthonormal
in-plane axes with the origin at \code{plane.point}. The axes are chosen automatically, so use
\code{point} if you need the position in model coordinates and \code{u}, \code{v} only for
distances within the plane (the example prints the in-plane radius
$\sqrt{u^2+v^2} = 1$ of the helix point $(0, 1, 1.25)$):

\begin{lstlisting}[style=text]
helix/plane z = 1.25: t = 1.250000 (expected 1.25), point (+0.000000, +1.000000, +1.250000), in-plane radius 1.000000
signed distance of the origin from the plane: -1.250
\end{lstlisting}

The helix is a C2 quintic approximation of the true helix within $10^{-7}$
(\cref{ch:curves}), so its parameter is ``turns'' only to that accuracy.

\section{Surface/surface intersections}\label{sec:isect-ss}

\code{intersect(s1, s2, options=SurfaceIntersectionOptions())} computes the intersection
\emph{curves} of two surfaces together with isolated touching points. It returns a
\apiref{SurfaceIntersection}:

\begin{center}
\begin{tabularx}{\linewidth}{@{}lX@{}}
\toprule
Member & Meaning \\
\midrule
\code{curves} & list of \apiref{IntersectionCurve} \\
\code{points} & list of \apiref{SurfaceSurfacePoint}: isolated points where the surfaces only
touch \\
\midrule
\code{IntersectionCurve.curve} & the intersection curve in 3D (a \code{NurbsCurve}) \\
\code{IntersectionCurve.pcurve1} & the same curve in the parameter plane of \code{s1}:
points $(u_1, v_1, 0)$, \emph{same parameter} as \code{curve} \\
\code{IntersectionCurve.pcurve2} & the same in the parameter plane of \code{s2} \\
\code{IntersectionCurve.closed} & \code{True} for a closed loop \\
\code{IntersectionCurve.points} & the traced points (\code{SurfaceSurfacePoint}) from which the
curves were fitted \\
\midrule
\code{SurfaceSurfacePoint.point} & 3D point \\
\code{SurfaceSurfacePoint.u1}, \code{v1}, \code{u2}, \code{v2} & parameters on both surfaces \\
\bottomrule
\end{tabularx}
\end{center}

The three curves of an \code{IntersectionCurve} share their parameter, which is what makes
them useful for trimming: for every $t$, $s_1(\mathrm{pcurve1}(t))$,
$s_2(\mathrm{pcurve2}(t))$ and $\mathrm{curve}(t)$ are the same point within the fit
tolerance. A p-curve is a \code{NurbsCurve} with $z=0$ whose $x$ and $y$ coordinates are the
surface parameters $u$ and $v$; this is the same convention that trim loops use
(\cref{ch:trimming}).

\subsection{Options}

\begin{center}
\begin{tabularx}{\linewidth}{@{}llX@{}}
\toprule
Option & Default & Meaning \\
\midrule
\code{tolerance} & \code{Tolerance()} & coincidence tolerances (seeds, ends on the domain
boundary, merging) \\
\code{chord_tolerance} & $0$ & maximal deviation of the traced polyline from the true curve
in model units; controls the marching step. $0$ means $10^{-4}$ of the model size. \\
\code{fit_tolerance} & $0$ & accuracy of the fitted curves: $|S_i(\mathrm{pcurve}_i(t)) -
\mathrm{curve}(t)|$ is verified between all fit points and exact points are inserted where
needed. $0$ means $10^{-7}$ of the model size. \\
\bottomrule
\end{tabularx}
\end{center}

The options are set by keyword in the constructor or as attributes afterwards:

\begin{lstlisting}
opts = nf.SurfaceIntersectionOptions(chord_tolerance=1e-6, fit_tolerance=1e-10)
opts.tolerance = nf.Tolerance(length=1e-8)
res = nf.intersect(s1, s2, opts)
\end{lstlisting}

``Model size'' is the extent of the two surfaces' control nets, so the defaults scale with
the model: a part measured in micrometres and the same part in metres give the same relative
accuracy.

\subsection{Example: sphere and cylinder}

A sphere of radius 2 and a cylinder of radius 1 about the same axis intersect in two circles
of radius 1 at $z=\pm\sqrt{3}$:

\example{09_surface_intersections.py}

The output shows how the circles come back:

\begin{lstlisting}[style=text]
sphere/cylinder: 4 curves, 0 touching points
  z = -1.732051, length 4.712389, closed=False, 119 traced points, p-curve residual 8.2e-09
  z = -1.732051, length 1.570796, closed=False, 39 traced points, p-curve residual 6.9e-08
  z = +1.732051, length 4.712389, closed=False, 119 traced points, p-curve residual 8.2e-09
  z = +1.732051, length 1.570796, closed=False, 39 traced points, p-curve residual 6.9e-08
total length 12.5663705982, exact 2 * 2 pi = 12.5663706144
chord 1e-6: 8910 traced points (default: 316)
\end{lstlisting}

Each circle consists of two open pieces of lengths $3\pi/2$ and $\pi/2$. The reason is the
\emph{seam}: the sphere and the cylinder are closed surfaces whose parameter domains have a
start and an end ($u=0$ and $u=1$) at the same 3D location. Intersection curves are split
where they cross the boundary of either parameter domain, because a p-curve cannot jump from
$u=1$ back to $u=0$. The two seams lie at different angles, hence two pieces per circle. If
you need the full circle, chain the pieces by their end points or use
\apiref{join} (\cref{ch:curves}). The residuals of the p-curves (below $10^{-7}$) are within
the default fit tolerance of $10^{-7}$ of the model size; the total length agrees with
$4\pi$ to $1.6\times10^{-8}$.

Tightening the chord tolerance from the default ($10^{-4}$ of the model size) to $10^{-6}$
increases the number of traced points from 316 to 8910. The chord tolerance mainly affects
robustness (small loops and sharp turns are not skipped) and run time; the accuracy of the
fitted curves is governed by \code{fit_tolerance}.

\subsection{Tangential contacts}

Where two surfaces only touch, the intersection is not transversal and cannot be traced by
marching. Isolated touching points are returned in \code{points}. Surfaces touching along a
whole curve are not traced as curves either; for two cylinders lying side by side the example
reports

\begin{lstlisting}[style=text]
touching cylinders: 0 curves, 131 points
\end{lstlisting}

i.e. the contact line is represented by points only. The Boolean operations treat such
contacts specially (\cref{sec:boolean-algorithm}); if you use the raw intersection,
check \code{points} as well as \code{curves}.

\subsection{Algorithm}

The surface/surface intersection proceeds in seven steps (details in the developer
documentation \code{docs/math/nurbs.md}):

\begin{enumerate}
\item \textbf{Seeds.} Pairs of Bézier patches with overlapping boxes are subdivided. When
their normal cones are separated, the pair cannot contain a closed loop (Sederberg's
criterion), so every intersection curve crosses the boundary of the pair: the boundary
iso-curves of each patch are intersected with the other patch by the curve/surface
intersection. Tangential regions stop at a size limit and are seeded by a minimum-norm
Gauss--Newton iteration.
\item \textbf{Marching} along $n_1\times n_2$ (limit normals at poles) with a predictor in
parameter space and a corrector onto both surfaces. The step size is controlled by the chord
tolerance in 3D \emph{and} in both parameter spaces. Branches end exactly on the domain
boundary, close into loops, or stop at tangential points.
\item \textbf{Knot lines.} Points are inserted exactly where a step crosses a knot line of
either surface, so that the fitted curves follow the C1 breaks of, e.g., rational circles.
\item \textbf{Verified fit.} The curve and both p-curves are interpolated piecewise; the
deviation is checked between all fit points against \code{fit_tolerance} and exact points
are added where needed (near poles the p-curves vary quickly).
\item \textbf{Duplicate seeds} on an already traced curve are skipped (compared in 3D and in
weighted parameter distance, so seams and poles are handled).
\item \textbf{Tangential contacts}: very short curves with parallel normals are reported as
points; curves with nearly parallel normals everywhere are tested for crossing versus
touching.
\item \textbf{Domain boundaries}: curves that touch a domain edge tangentially are continued
from another seed, and open curves ending at coinciding boundary points are joined.
\end{enumerate}

\section{Planar sections of surfaces}\label{sec:planar-section}

\apiref{planar_section} intersects one surface with an infinite plane:

\begin{lstlisting}
res = nf.planar_section(surface, nf.Plane(point, normal), options=nf.SurfaceIntersectionOptions())
\end{lstlisting}

Internally the plane is represented by a rectangle that covers the projection of the
surface's control points (plus a margin), so no section curve is missed, and the ordinary
surface/surface intersection is used. The result is a \code{SurfaceIntersection} like above,
with one difference: \code{pcurve2} (and \code{u2}, \code{v2} of the points) are expressed
in the plane's own orthonormal coordinates $(s,t)$ with the origin at \code{plane.point} --
the same in-plane coordinates as for \code{intersect(curve, plane)}. \code{pcurve1} is the
p-curve on the surface, as usual.

\apiref{planar_sections} slices a surface with \code{count} parallel planes,
\code{plane.point + k * spacing * normal} for $k = 0,\dots,\mathrm{count}-1$, and returns
one \code{SurfaceIntersection} per plane:

\begin{lstlisting}[style=text]
section z = 1: 1 curve(s), length 10.88279617 (2 pi sqrt 3 = 10.88279619), first point in plane coords radius 1.73205081
z = -1.5: 8.31187287 (exact 8.31187288)
z = -0.5: 12.16733602 (exact 12.16733603)
z = +0.5: 12.16733602 (exact 12.16733603)
z = +1.5: 8.31187287 (exact 8.31187288)
\end{lstlisting}

\begin{note}[Surfaces versus solids]
\code{planar_section} works on a single, untrimmed surface and returns unordered curve
pieces. To section a \emph{solid} -- with trimmed faces, closed and oriented loops, holes and
islands -- use \apiref{section} on a \code{Shell}, described in \cref{ch:slicing}.
\end{note}

\section{Closest points}\label{sec:closest-point}

Point projection (``inversion'') answers the question: which parameter of a curve or surface
lies closest to a given point? \NF{} provides a global and a local variant for both curves
and surfaces.

\begin{center}
\begin{tabularx}{\linewidth}{@{}lX@{}}
\toprule
Call & Result \\
\midrule
\code{closest_point(curve, p, tolerance=Tolerance())} & \apiref{CurvePoint}: \code{u},
\code{point}, \code{distance} \\
\code{closest_point(surface, p, tolerance=Tolerance())} & \apiref{SurfacePoint}: \code{u},
\code{v}, \code{point}, \code{distance} \\
\code{closest_point_local(curve, p, u0, tolerance)} & \code{CurvePoint}, Newton iteration
from \code{u0} \\
\code{closest_point_local(surface, p, u0, v0, tolerance)} & \code{SurfacePoint}, Newton
iteration from $(u_0,v_0)$ \\
\bottomrule
\end{tabularx}
\end{center}

The global variant subdivides the object into Bézier pieces, discards pieces whose bounding
box is farther away than the best candidate so far, and refines the remaining candidates by
Newton's method; it always returns the global minimum of the distance (up to ties). The
local variant only runs the Newton iteration from the given start parameter and is much
faster; use it when you already know a good starting point, e.g. while tracking a moving
point along a curve.

\example{09_closest_points.py}

\begin{lstlisting}[style=text]
circle: u = 0.896447, point (1.200000, 1.600000), distance 3.000000 (expected 5 - 2 = 3)
line, local from u0 = 0.5: u = 0.750000, distance 1.000000
sphere: (u, v) = (0.250000, 0.750000), distance 3.242641 (expected 3 sqrt 2 - 1 = 3.242641)
on the axis: distance 4.000000, point z = 1.000000
\end{lstlisting}

When the minimum is not unique -- a point on the axis of a sphere is equally close to a
whole parallel circle, the centre of a circle to every point of it -- one of the minimisers
is returned. The distance is correct in any case.

\section{Performance and robustness notes}\label{sec:isect-notes}

\begin{itemize}
\item \textbf{Degree and knot count.} The subdivision works on Bézier pieces; curves and
surfaces with many knot spans produce many pieces. Run time scales roughly with the number of
overlapping piece pairs. Simplify inputs with \code{remove_knot} or approximate them if they
are heavily over-refined (\cref{ch:curves}).
\item \textbf{Parallel seeds.} The seed search of the surface/surface intersection is
parallelised internally; the rest of the algorithm runs on one thread. Several independent
intersections can be run in parallel from Python threads, since all objects are immutable.
\item \textbf{Units.} All defaults that depend on the ``model size'' scale with the inputs;
the absolute \code{Tolerance} fields do not. Very large models (kilometres) or very small ones
(nanometres) may need an explicit \code{Tolerance(length=...)}.
\item \textbf{Errors.} Invalid inputs raise \code{GeometryError}. An empty list, or a
result without curves and points, means that nothing was found within the tolerance.
\end{itemize}
